\documentclass{article}
\usepackage[T1]{fontenc}
\usepackage{lmodern}
\usepackage{graphicx}
\usepackage{amsmath,amssymb}
\usepackage[a4paper,margin=2cm]{geometry}
\usepackage[absolute,overlay]{textpos}
\setlength{\TPHorizModule}{1mm}
\setlength{\TPVertModule}{1mm}
\usepackage[indonesian]{babel}
\usepackage{hyperref}
\setlength{\emergencystretch}{2em}
% unit: Halaman 88

\begin{document}

% === Halaman 88 (python/native) ===

3. Cocok untuk perangkat dengan sumber daya terbatas

Karena operasi internalnya sederhana, Salsa20 dapat diimplementasikan dengan 

baik pada berbagai perangkat, termasuk perangkat yang kemampuan komputasinya 

terbatas.

4. Cocok untuk implementasi perangkat lunak

Salsa20 dapat memperoleh performa yang baik hanya dengan operasi integer 

dasar. Karena itu Salsa20 sangat populer dalam komunitas kriptografi, dan 

desainnya kemudian menjadi dasar bagi ChaCha, yang akhirnya lebih banyak 

digunakan dalam protokol modern.

5. Menghasilkan blok keystream yang besar

Satu blok Salsa20 menghasilkan 16 × 32 bit = 512 bit = 64 byte keystream. Jika 

plaintext panjang, counter tinggal dinaikkan untuk menghasilkan blok berikutnya.

88

\end{document}
